\begin{frame}{Cuerpos}

Trabajaremos con matrices cuyos coeficientes pertenecen a un \alert{cuerpo} \(K\).

\medskip

Para nosotros, un cuerpo es un conjunto de ``números'' en el que podemos:

\begin{itemize}
    \item sumar,
    \item restar,
    \item multiplicar,
    \item dividir por cualquier número distinto de cero.
\end{itemize}

\medskip

Ejemplos de cuerpos:
\[
\mathbb{Q},
\qquad
\mathbb{R},
\qquad
\mathbb{C}.
\]

\medskip

En cambio, \(\mathbb{N}\) y \(\mathbb{Z}\) no son cuerpos:

\begin{itemize}
    \item en \(\mathbb{N}\) no siempre podemos restar:
    \[
    2-3\notin\mathbb{N};
    \]
    \item en \(\mathbb{Z}\) no siempre podemos dividir por un número no nulo:
    \[
    \frac{1}{2}\notin\mathbb{Z}.
    \]
\end{itemize}

\end{frame}